\documentclass{article}
\input{../../../lib.sty}

\title{\texttt{fn new\_poisson\_rdp\_curve}}
\author{Michael Shoemate}

\begin{document}
\maketitle

Proves soundness of \texttt{fn new\_poisson\_rdp\_curve} in \texttt{mod.rs}.

\section{Hoare Triple}
\subsection*{Precondition}
\texttt{base\_curve} upper-bounds the base Rényi divergence at every order $> 1$, and
\texttt{mean} $\ge 0$.

\subsection*{Pseudocode}
\lstinputlisting[language=Python,firstline=2,escapechar=|]{./pseudocode/new_poisson_rdp_curve.py}

\subsection*{Postcondition}
\begin{theorem}
\label{postcondition}
For every $\alpha > 1$, the returned function either returns \texttt{Err(e)} due to a propagated
helper error, or returns an upper bound on the implemented Poisson expression
\[
\epsilon(\alpha)
\;+\;
\texttt{mean}\cdot \delta(\alpha)
\;+\;
\frac{\log(\texttt{mean})}{\alpha - 1},
\]
where $\epsilon(\cdot)$ is the base Rényi divergence curve and
\[
\delta(\alpha) = \alpha^{-1}\!\left(1 - \alpha^{-1}\right)^{\alpha - 1}.
\]
\end{theorem}

\section{Proof}
Fix $\alpha > 1$.

\begin{lemma}
\label{special-cases}
If \texttt{mean = 0}, the returned value is $0$, which equals the displayed expression.
If $\alpha \le 1$, the returned value is $+\infty$, which is a valid upper bound.
\end{lemma}

\begin{proof}
When \texttt{mean = 0}, both additive terms involving \texttt{mean} vanish.
For $\alpha \le 1$, $+\infty$ upper-bounds any real value.
\end{proof}

\begin{lemma}
\label{admissibility}
If the admissibility check on line \ref{line:admissibility} passes, then the exact base Rényi
divergence satisfies the Theorem 6 admissibility condition.
\end{lemma}

\begin{proof}
Let $\epsilon(\alpha)$ denote the exact base Rényi divergence curve.
By precondition, \texttt{eps\_alpha} $\ge \epsilon(\alpha)$.
By the proof of \texttt{log\_alpha\_over\_alpha\_minus\_one\_lo},
\[
\texttt{max\_eps\_lo} \le \log\!\left(\frac{\alpha}{\alpha - 1}\right).
\]
Therefore \texttt{eps\_alpha} $\le$ \texttt{max\_eps\_lo} implies
\[
\epsilon(\alpha) \le \log\!\left(\frac{\alpha}{\alpha - 1}\right),
\]
which is exactly the admissibility condition.
\end{proof}

\begin{lemma}
\label{admissibility-fails}
If the admissibility check on line \ref{line:admissibility} fails, the function returns
$+\infty$, which is a valid upper bound on the displayed expression.
\end{lemma}

\begin{proof}
This is exactly the branch taken by the pseudocode, and $+\infty$ upper-bounds any real value.
\end{proof}

\begin{lemma}
\label{delta}
If \texttt{mean > 0}, $\alpha > 1$, and the admissibility check passes, then the value
\texttt{delta\_hi} computed by the function satisfies
\[
\texttt{delta\_hi} \ge \delta(\alpha).
\]
\end{lemma}

\begin{proof}
We have \texttt{inv\_alpha\_hi} $\ge 1/\alpha$ and \texttt{inv\_alpha\_lo} $\le 1/\alpha$.
Hence
\[
\texttt{one\_minus\_inv\_alpha\_hi} \ge 1 - \frac{1}{\alpha},
\]
so after subtracting $1$,
\[
\texttt{one\_minus\_inv\_alpha\_minus\_one\_hi} \ge -\frac{1}{\alpha}.
\]
Both values lie in $(-1,0)$, so upward-rounded \texttt{ln\_1p} gives
\[
\texttt{ln\_one\_minus\_inv\_alpha\_hi}
\ge
\ln\!\left(1 - \frac{1}{\alpha}\right).
\]
Also \texttt{alpha\_minus\_one\_lo} $\le \alpha - 1$ with \texttt{alpha\_minus\_one\_lo} $> 0$.
Since the logarithm above is negative, multiplying first by the smaller positive factor and then
by the larger negative factor is conservative:
\[
\texttt{exponent\_hi}
\ge
(\alpha - 1)\ln\!\left(1 - \frac{1}{\alpha}\right).
\]
Upward-rounded exponentiation yields
\[
\texttt{power\_hi}
\ge
\left(1 - \frac{1}{\alpha}\right)^{\alpha - 1}.
\]
Multiplying by \texttt{inv\_alpha\_hi} $\ge 1/\alpha$ gives the claim.
\end{proof}

\begin{lemma}
\label{main}
If \texttt{mean > 0}, $\alpha > 1$, the admissibility check passes, and the function returns
\texttt{Ok(out)}, then \texttt{out} upper-bounds the displayed Poisson expression.
\end{lemma}

\begin{proof}
By precondition, \texttt{eps\_alpha} $\ge \epsilon(\alpha)$.
By Lemma~\ref{delta},
\[
\texttt{term1\_hi} \ge \texttt{mean}\cdot \delta(\alpha).
\]
By the proof of \texttt{log\_mean\_hi},
\[
\texttt{log\_mean\_hi\_} \ge \log(\texttt{mean}).
\]
Since \texttt{alpha\_minus\_one\_lo} $\le \alpha - 1$ with a positive denominator,
\[
\texttt{term2\_hi} \ge \frac{\log(\texttt{mean})}{\alpha - 1}.
\]
Finally, \texttt{out} is formed by upward-rounded addition of these three upper bounds.
\end{proof}

\begin{proof}[Proof of Theorem~\ref{postcondition}]
By Lemma~\ref{special-cases}, Lemma~\ref{admissibility-fails}, Lemma~\ref{admissibility}, and
Lemma~\ref{main}.
\end{proof}

\end{document}
